\documentclass{article}

\usepackage{amsmath,amssymb}
\usepackage{xurl}
\usepackage[hidelinks]{hyperref}
\setlength{\emergencystretch}{2em}

\input{command}

\title{The Rotation Sign in Wolf's Spinor Exponential Formula}
\author{}
\date{2026-08-24}

\begin{document}
\maketitle
\gapnote{false-source}{resolved}

\section{The printed formula}

Let $\sigma_1,\sigma_2,\sigma_3$ be the standard Pauli matrices. Wolf defines
the spatial rotation generators by
\begin{align}
    R_i
        &=\sum_{j,k=1}^{3}
          \varepsilon_{ijk}\ketbra{k}{j}.
    \label{eq:wolf_ch2_spinor_rotation_sign_rotation_generators}
\end{align}
For $\vec n\in\mathbb R^3$, Equation~(2.44) of
Ref.~\cite{Wolf2012Quantum} prints the correspondence
\begin{align}
    U
        &=\exp\!\left(-\frac{i}{2}\vec n\mathbin{\cdot}\vec\sigma\right),
    \label{eq:wolf_ch2_spinor_rotation_sign_printed_spinor}\\
    U
        &\longmapsto
          \exp\!\left(-\vec n\mathbin{\cdot}\vec R\right).
    \label{eq:wolf_ch2_spinor_rotation_sign_printed_rotation}
\end{align}
The local source is
\path{Notes/WolfNoteTexSource/ch02_representations.tex}, lines~1070--1077.

\section{The sign inconsistency}

Take
\begin{align}
    \vec n
        &=(0,0,\theta).
    \label{eq:wolf_ch2_spinor_rotation_sign_z_axis_parameter}
\end{align}
With the standard matrix
\begin{align}
    \sigma_2
        &=\begin{pmatrix}
            0 & -i\\
            i & 0
          \end{pmatrix}
    \label{eq:wolf_ch2_spinor_rotation_sign_pauli_two}
\end{align}
and the column-vector convention used in the source, conjugation gives
\begin{align}
    e^{-i\theta\sigma_3/2}\sigma_1e^{i\theta\sigma_3/2}
        &=\cos\theta\,\sigma_1+\sin\theta\,\sigma_2.
    \label{eq:wolf_ch2_spinor_rotation_sign_pauli_conjugation}
\end{align}
Thus the Pauli-coordinate vector of $\sigma_1$ rotates toward the positive
second coordinate.

By
(\ref{eq:wolf_ch2_spinor_rotation_sign_rotation_generators}), the third
generator satisfies
\begin{align}
    R_3\ket{1}
        &=\ket{2}.
    \label{eq:wolf_ch2_spinor_rotation_sign_generator_action}
\end{align}
Therefore
\begin{align}
    e^{+\theta R_3}\ket{1}
        &=\cos\theta\,\ket{1}+\sin\theta\,\ket{2},
    \label{eq:wolf_ch2_spinor_rotation_sign_positive_rotation_action}\\
    e^{-\theta R_3}\ket{1}
        &=\cos\theta\,\ket{1}-\sin\theta\,\ket{2}.
    \label{eq:wolf_ch2_spinor_rotation_sign_negative_rotation_action}
\end{align}
The conjugation in
(\ref{eq:wolf_ch2_spinor_rotation_sign_pauli_conjugation}) agrees with
(\ref{eq:wolf_ch2_spinor_rotation_sign_positive_rotation_action}), not with
(\ref{eq:wolf_ch2_spinor_rotation_sign_negative_rotation_action}). Hence the
two negative signs in the printed correspondence cannot both be correct.

\section{Correction and formalization}

For the spinor in
(\ref{eq:wolf_ch2_spinor_rotation_sign_printed_spinor}) and the generators in
(\ref{eq:wolf_ch2_spinor_rotation_sign_rotation_generators}), the corrected
congruence action is
\begin{align}
    \exp\!\left(-\frac{i}{2}\vec n\mathbin{\cdot}\vec\sigma\right)
        &\longmapsto
          \exp\!\left(+\vec n\mathbin{\cdot}\vec R\right).
    \label{eq:wolf_ch2_spinor_rotation_sign_corrected_correspondence}
\end{align}
Equivalently, the negative Lorentz exponential printed in Equation~(2.44) of
Ref.~\cite{Wolf2012Quantum} corresponds to
\begin{align}
    U
        &=\exp\!\left(+\frac{i}{2}\vec n\mathbin{\cdot}\vec\sigma\right).
    \label{eq:wolf_ch2_spinor_rotation_sign_equivalent_spinor}
\end{align}

The corrected relation
(\ref{eq:wolf_ch2_spinor_rotation_sign_corrected_correspondence}) is
formalized directly. The formal development also contains the closed spinor
exponential, Rodrigues' formula, and the consistent positive-angle
diagonal-$\mathrm{SU}(2)$ calculation.\footnote{Formal declarations:
\leanid{Wolf.spinorMatrix\_rotationExpSL2},
\leanid{Wolf.exp\_neg\_I\_smul\_pauliVector}, and
\leanid{Wolf.exp\_smul\_rotationGenerator} in
\doclink{QICLean/Channel/LorentzNormalForm/SpinorExponential}; the
positive-angle diagonal calculation is
\leanid{SpinCover.R\_su2Diag\_eq\_rotZ}.}

\section{Verdict}

\emph{False source.} Equation~(2.44) of Ref.~\cite{Wolf2012Quantum} contains a
genuine sign error. This is not a convention change introduced by the
formalization: the source explicitly fixes the Pauli matrices, the generators
in (\ref{eq:wolf_ch2_spinor_rotation_sign_rotation_generators}), and the
column-vector action. The boost part of Equation~(2.44) is unaffected.

\bibliographystyle{alpha}
\bibliography{references}

\end{document}
